% s-225-p225-m001; original PDF 225, 226
\Aphorism{s-225-p225-m001}{317}
Ich habe nie eine philosophischere Gebärde gesehen als jene Umarmung, mit der Mitterwurzer als Hamlet von Ophelia Abschied nimmt — wobei es gleichgültig ist, ob der Darsteller die Geberde oder ihre Deutung oder keins von beiden beabsichtigte. Es handelt sich um das Gespräch, nach dem Ophelia die Zerstörung des edelsten Geistes beklagt; sie hat soeben ein Indicium dafür bekommen, an Hamlets Umarmung, die er ausführt, ohne das Weib körperlich zu berühren. Er gleitet mit einer suchenden, saugenden Tastbewegung an ihrem „Ätherleib“ hin, längs einer unsichtbaren, atmosphärischen Hülle, die ihm, wie es scheint, eine völlige Annäherung verwehrt. Das ist ein eminent physiologisches Symbol Hamlets, des seelischen Weichthiers, des Gallertmenschen, der durch Bewusstsein und Metaphysik seiner Persönlichkeitshülle entkleidet, hyperästhetisch und berührungsscheu geworden ist. Für ihn, den Enthäuteten, mit allen Nerven Blossgelegten, härtet und verdickt sich die Haut der Objecte über die Grenzen ihres realen Umfangs hinaus; ihre Ausstrahlung schon, die auf ihrer Oberfläche niedergeschlagene Gasschicht, leistet ihm körperlichen Widerstand — ich drücke mich physicalisch aus, weil für solche psychischen Thatbestände noch keine eigene Formelsprache erfunden ist. Wie anders steht der active, instinctive Mensch zu den Dingen! „Mit einem Griff zergreif’ ich den Quark!“ Der hat harte Hände und lebt in einer starken subjectiven Hülle von gewaltsam selbstherrlicher Ungerechtigkeit; ihm erweicht und verflüssigt sich die Umwelt, die sich dem objectiven Bewusstseinsmenschen verkantet und versteinert. Hamlet nimmt ungepanzert den Kampf mit dem Leben auf und unterliegt; Siegfried, der Gehörnte, Gehäutete, ist des Sieges gewiss. Siegfried ringt Brünnhilden nieder, Hamlet lässt Ophelien unberührt los.
\par

